\section{RQ2: Context-Insensitive Metrics for Hallucination Detection}
\label{sec:rq2_statistical_methods}

This section documents the methodology for Research Question 2: \textit{Can context-insensitive (CI) metrics computed from generator LLM logits be used to detect hallucinations?}

\subsection{CI Metrics: Mathematical Definitions}
\label{sec:ci_metrics_definitions}

Context-insensitive metrics are computed from the generator LLM's token-level log-probabilities during edge generation. These metrics capture model uncertainty without requiring external context or reasoning.

\subsubsection{Generator Perplexity}

Perplexity measures the exponential of average negative log-likelihood over generated tokens, quantifying how ``surprised'' the model is by its own output.

\begin{equation}
\text{Perplexity} = \exp\left(\frac{1}{n}\sum_{i=1}^{n} -\log P(x_i)\right)
\label{eq:perplexity}
\end{equation}

where:
\begin{itemize}
    \item $n$ = number of generated tokens in the causal motivation
    \item $P(x_i)$ = probability of token $x_i$ (from LLM logprobs: $P(x_i) = \exp(\text{logprob}_i)$)
\end{itemize}

\textbf{Implementation:} \texttt{logit\_metrics.py:compute\_perplexity()}. Average NLL capped at 10.0 to prevent overflow (max perplexity $\approx 22{,}026$).

\textbf{Interpretation:} Higher perplexity $\rightarrow$ greater model uncertainty $\rightarrow$ higher hallucination probability.

\subsubsection{Generator Minimum Probability}

Minimum probability identifies the ``weakest link'' in the generated sequence---the token about which the model was least confident.

\begin{equation}
\text{MinProb} = \min_{i \in \{1,\ldots,n\}} P(x_i)
\label{eq:minprob}
\end{equation}

where $P(x_i) = \exp(\text{logprob}_i)$ from LLM logprobs.

\textbf{Implementation:} \texttt{logit\_metrics.py:compute\_min\_prob()}.

\textbf{Interpretation:} Lower minimum probability $\rightarrow$ presence of highly uncertain token $\rightarrow$ higher hallucination probability.

\subsubsection{Generator Maximum Window Entropy}

Maximum window entropy detects localized regions of high uncertainty by computing Shannon entropy over sliding windows of top-$k$ token distributions.

\begin{equation}
\text{MaxWindowEntropy} = \max_{w \in \mathcal{W}} \left( \frac{1}{|w|} \sum_{t \in w} H_k(t) \right)
\label{eq:maxent}
\end{equation}

where the local entropy at position $t$ over the top-$k$ tokens is:

\begin{equation}
H_k(t) = -\sum_{j=1}^{k} \tilde{P}_j(t) \cdot \ln \tilde{P}_j(t)
\label{eq:local_entropy}
\end{equation}

and $\tilde{P}_j(t)$ is the normalized probability over top-$k$ alternatives:

\begin{equation}
\tilde{P}_j(t) = \frac{P_j(t)}{\sum_{m=1}^{k} P_m(t)}
\label{eq:normalized_prob}
\end{equation}

\textbf{Parameters:} $k = 5$ (top-$k$ tokens), window size $= 5$ tokens, sliding with stride 1.

\textbf{Implementation:} \texttt{logit\_metrics.py:compute\_max\_window\_entropy()}. Uses \texttt{np.convolve} for rolling average.

\textbf{Interpretation:} Higher max entropy $\rightarrow$ region where model was uncertain among alternatives $\rightarrow$ higher hallucination probability.

\subsubsection{Generator Cosine Similarity}

Cosine similarity measures semantic alignment between the generated causal motivation and fetched citation text chunks.

\begin{equation}
\text{CosineSim} = \max_{c \in \mathcal{C}} \frac{\mathbf{e}_{\text{narr}} \cdot \mathbf{e}_c}{\|\mathbf{e}_{\text{narr}}\| \cdot \|\mathbf{e}_c\|}
\label{eq:cosine}
\end{equation}

where:
\begin{itemize}
    \item $\mathbf{e}_{\text{narr}} \in \mathbb{R}^{d}$ = embedding of generated causal narrative (motivation)
    \item $\mathbf{e}_c \in \mathbb{R}^{d}$ = embedding of citation chunk $c$
    \item $\mathcal{C}$ = set of all citation chunks (up to 5 articles, chunked with 20\% overlap)
    \item $d = 1536$ for \texttt{text-embedding-3-small} (OpenAI)
\end{itemize}

\textbf{Implementation:} \texttt{logit\_metrics.py:compute\_alignment\_batch()}. Chunking: $\sim$2000 characters per chunk, 20\% overlap.

\textbf{Interpretation:} Higher cosine similarity was \textit{expected} to indicate better grounding (lower hallucination). \textbf{Counterintuitive finding:} Higher similarity correlates with \textit{more} hallucinations (see Section~\ref{sec:cosine_paradox}).

\subsubsection{Summary of CI Metrics}

\begin{table}[h]
\centering
\begin{tabular}{llcc}
\toprule
\textbf{Metric} & \textbf{Formula} & \textbf{Expected} & \textbf{Observed} \\
\midrule
Perplexity & $\exp\left(\frac{1}{n}\sum -\log P(x_i)\right)$ & $\uparrow \Rightarrow$ halluc & $\uparrow \Rightarrow$ halluc \\
Min Prob & $\min_i P(x_i)$ & $\downarrow \Rightarrow$ halluc & $\downarrow \Rightarrow$ halluc \\
Max Window Entropy & $\max_w \text{mean}(H_k(t))$ & $\uparrow \Rightarrow$ halluc & $\uparrow \Rightarrow$ halluc \\
Cosine Similarity & $\max_c \cos(\mathbf{e}_{\text{narr}}, \mathbf{e}_c)$ & $\downarrow \Rightarrow$ halluc & $\uparrow \Rightarrow$ halluc$^*$ \\
\bottomrule
\end{tabular}
\caption{CI metrics summary. $^*$Counterintuitive: higher similarity predicts \textit{more} hallucinations.}
\label{tab:ci_metrics_summary}
\end{table}

\subsection{Statistical Analysis Methods}

This section documents the statistical methodology employed to evaluate whether CI metrics predict hallucinations in LLM-generated causal discovery outputs. All analyses were implemented in Python 3.11.x using standard scientific computing libraries and employ rigorous validation with multiple comparison correction.

\subsection{Study Design and Data Structure}

\subsubsection{Experimental Data Structure}

The RQ2 analysis operates on \textbf{63 deduplicated experimental files} from two experiment types:

\begin{table}[h]
\centering
\begin{tabular}{lcccc}
\toprule
\textbf{Experiment Type} & \textbf{CLDs} & \textbf{Runs} & \textbf{Prompts} & \textbf{Files} \\
\midrule
Citation (GT Lit) & 3 & 3 & 4 & 36 \\
Correctness (GT Lit) & 3 & 3 & 3 & 27 \\
\midrule
\textbf{Total} & & & & \textbf{63} \\
\bottomrule
\end{tabular}
\caption{RQ2 experimental file structure. Each file contains one complete CLD generation.}
\label{tab:rq2_file_structure}
\end{table}

\paragraph{Causal Loop Diagrams (CLDs).}
\begin{itemize}
    \item \textbf{Depressive symptoms}: Response to stressor in healthy adults ($\sim$210 edges/file)
    \item \textbf{Social norms}: Social norms and obesity prevalence ($\sim$110 edges/file)
    \item \textbf{Emergency department}: Older persons ED visits and interactions ($\sim$1190 edges/file)
\end{itemize}

\paragraph{Prompt Types.}
\begin{itemize}
    \item \textbf{Citation experiment (4 prompts):} baseline, cot, mechanistic\_original, mechanistic\_lit
    \item \textbf{Correctness experiment (3 prompts):} baseline, cot, mechanistic
\end{itemize}

\paragraph{Runs.} Three independent runs (run\_1, run\_2, run\_3) per (CLD $\times$ prompt) combination, using different LLM generation seeds.

\paragraph{Total Sample Size.} After pooling all 63 files: $\sim$31,700 total edges. After removing rows with missing CI metrics: $\sim$16,500 edges (52.1\% retention). Hallucination rate: $\sim$15\% (2,514 hallucinations, 13,993 correct).

\subsubsection{Cross-CLD Validation Design}

The analysis employed a multi-dataset validation design to assess generalizability:

\begin{itemize}
    \item \textbf{Primary dataset structure:} 3 Causal Loop Diagrams (CLDs) from the health domain
    \begin{itemize}
        \item CLD 1: Depressive symptoms in response to a stressor (health psychology)
        \item CLD 2: Social norms and obesity prevalence (health behavior)
        \item CLD 3: Older persons emergency department visits (healthcare systems)
    \end{itemize}
    \item \textbf{Replication structure:} 3 independent runs per CLD $\times$ prompt combination
    \item \textbf{Total files:} 63 deduplicated experimental files
    \item \textbf{Meta-analytic approach:} One-sample t-tests aggregating performance metrics across files
\end{itemize}

\subsubsection{Ground Truth Definition}

\paragraph{Hallucination labeling.} For GT Lit experiments (citation and correctness), hallucinations are defined as:
\begin{equation}
\text{is\_hallucination} = \mathbb{I}(\text{Classification} = \text{FP}) \lor \mathbb{I}(\text{Classification} = \text{FN})
\label{eq:halluc_def}
\end{equation}

where:
\begin{itemize}
    \item \textbf{False Positive (FP):} Edge generated by LLM but absent from validation CLD
    \item \textbf{False Negative (FN):} Edge present in validation CLD but not generated by LLM
\end{itemize}

\paragraph{Data availability.} CI metrics (perplexity, min\_prob, max\_window\_entropy, cosine\_similarity) computed only for edges with LLM-generated motivation text. For FN edges (missed by generator), no generation occurred, hence no logprobs available---these are excluded from individual CI metric analysis but counted in classification performance.

\paragraph{Note on corruption experiments.} GT Synth (corruption) experiments are excluded from RQ2 analysis because corruption is inserted \textit{post-generation}. Generator CI metrics (computed during generation) cannot retroactively detect corruption inserted afterwards.

\subsection{Statistical Test Procedures}

\subsubsection{Correlation Analysis}

\paragraph{Pearson Product-Moment Correlation.} For metric values \(X = (x_1, \ldots, x_n)\) and binary hallucination labels \(Y = (y_1, \ldots, y_n)\):
\[
r = \frac{\sum_{i=1}^{n}(x_i - \bar{x})(y_i - \bar{y})}{\sqrt{\sum_{i=1}^{n}(x_i - \bar{x})^2 \sum_{i=1}^{n}(y_i - \bar{y})^2}}
\]
where \(\bar{x} = \frac{1}{n}\sum_{i=1}^{n}x_i\) and \(\bar{y} = \frac{1}{n}\sum_{i=1}^{n}y_i\).

\textbf{Hypothesis test:} Two-tailed, \(H_0: \rho = 0\), \(\alpha = 0.05\) \\
\textbf{Assumptions:} Bivariate normality (relaxed for large \(n\) via CLT)

\paragraph{Spearman Rank Correlation.} Non-parametric alternative robust to non-linearity and outliers:
\[
\rho_s = 1 - \frac{6\sum_{i=1}^{n}d_i^2}{n(n^2-1)}
\]
where \(d_i = \text{rank}(x_i) - \text{rank}(y_i)\).

\textbf{Use case:} Validate Pearson results when distributional assumptions violated

\subsubsection{Classification Performance Analysis}

\paragraph{Receiver Operating Characteristic (ROC) Analysis.} For each threshold \(t\) applied to CI metric \(X\):
\[
\text{TPR}(t) = \frac{\text{TP}(t)}{\text{TP}(t) + \text{FN}(t)}, \quad \text{FPR}(t) = \frac{\text{FP}(t)}{\text{FP}(t) + \text{TN}(t)}
\]

Area under ROC curve quantifies overall discriminative ability:
\[
\text{AUC} = \int_0^1 \text{TPR}(\text{FPR}^{-1}(x))\,dx
\]

\textbf{Implementation:} Standard ROC curve algorithm computes (FPR, TPR) pairs at all unique threshold values; AUC computed via trapezoidal rule integration.

\textbf{Interpretation scale:}
\begin{itemize}
    \item AUC > 0.9: Excellent discrimination
    \item 0.8--0.9: Good
    \item 0.7--0.8: Acceptable
    \item 0.6--0.7: Poor
    \item < 0.6: Fail (no better than chance)
    \item AUC = 0.5: Chance performance (random classifier)
\end{itemize}

\paragraph{Optimal Threshold Selection.} Thresholds selected to maximize F\(_1\) score:
\[
F_1 = \frac{2 \cdot \text{precision} \cdot \text{recall}}{\text{precision} + \text{recall}} = \frac{2\text{TP}}{2\text{TP} + \text{FP} + \text{FN}}
\]

\textbf{Rationale:} F\(_1\) balances precision (minimize false alarms) against recall (minimize missed hallucinations), appropriate for imbalanced classes.

\textbf{Implementation:} Grid search over percentile-based thresholds:
\begin{itemize}
    \item Candidate thresholds: \(t \in \{\text{percentile}_{10}, \text{percentile}_{15}, \ldots, \text{percentile}_{90}\}\) of metric distribution
    \item For each threshold \(t\), compute F\(_1\)(t) based on predicted labels
    \item Select optimal: \(t^* = \arg\max_t F_1(t)\)
\end{itemize}

This percentile-based grid search ensures coverage across the metric's empirical distribution while avoiding extreme values that may lead to degenerate classifications (e.g., all positive or all negative predictions).

\subsubsection{Permutation Tests}

\textbf{Purpose:} Permutation testing provides \textit{non-parametric} validation of correlation significance without assuming a specific distribution (e.g., normality). This addresses the limitation that Pearson's \(p\)-values assume bivariate normality, which may be violated for CI metrics.

\textbf{Null Hypothesis:} \(H_0\): metric values and hallucination labels are independent (no association).

\paragraph{Algorithm.}
\begin{enumerate}
    \item Compute observed correlation: \(r_{\text{obs}} = \text{corr}(X, Y)\)
    \item For \(i = 1, \ldots, n_{\text{perm}}\) where \(n_{\text{perm}} = 10{,}000\):
    \begin{enumerate}
        \item Randomly permute labels: \(Y^{(i)} = \text{permute}(Y)\)
        \item Compute permuted correlation: \(r_{\text{perm},i} = \text{corr}(X, Y^{(i)})\)
    \end{enumerate}
    \item Compute two-tailed p-value:
    \[
    p = \frac{1}{n_{\text{perm}}}\sum_{i=1}^{n_{\text{perm}}} \mathbb{I}(|r_{\text{perm},i}| \geq |r_{\text{obs}}|)
    \]
\end{enumerate}

where \(\mathbb{I}(\cdot)\) is the indicator function.

\textbf{Null distribution:} Under \(H_0\), permuted correlations form empirical null distribution centered at zero.

\textbf{Advantages over parametric tests:}
\begin{itemize}
    \item No distributional assumptions
    \item Exact Type I error control (conditional on data)
    \item Robust to outliers and skewness
\end{itemize}

\paragraph{Obtained Results.} Single-CLD analysis (Depressive Symptoms, n=182):
\begin{itemize}
    \item \textbf{perplexity}: \(r_{\text{obs}} = +0.331\), \(p_{\text{perm}} < 0.0001\) (10,000/10,000 permutations had \(|r| < |r_{\text{obs}}|\))
    \item \textbf{max\_window\_entropy}: \(r_{\text{obs}} = +0.327\), \(p_{\text{perm}} < 0.0001\)
    \item \textbf{min\_prob}: \(r_{\text{obs}} = -0.175\), \(p_{\text{perm}} = 0.0144\)
    \item \textbf{cosine\_similarity}: \(r_{\text{obs}} = +0.293\), \(p_{\text{perm}} < 0.0001\)
\end{itemize}

All generator CI metrics demonstrated p-values below \(\alpha = 0.05\), confirming associations not attributable to chance (\texttt{rq2\_permutation\_tests.csv}).

\paragraph{Multiple Comparison Correction.} Testing \(k = 4\) generator CI metrics inflates family-wise error rate (FWER). We applied Holm-Bonferroni sequential correction:

For ordered p-values \(p_{(1)} \leq p_{(2)} \leq \cdots \leq p_{(k)}\), reject \(H_{0,(i)}\) if:
\[
p_{(i)} \leq \frac{\alpha}{k - i + 1}
\]

\textbf{Implementation:} Sequential Holm-Bonferroni method with \(\alpha = 0.05\)

\textbf{FWER control:} Guarantees \(P(\text{any false rejection}) \leq \alpha\) under global null.

\textbf{Correction results:} All 4 generator CI metrics remained significant after Holm-Bonferroni correction, confirming robust family-wise significance control:
\begin{itemize}
    \item perplexity: \(p_{\text{corrected}} < 0.0001\)
    \item max\_window\_entropy: \(p_{\text{corrected}} < 0.0001\)
    \item cosine\_similarity: \(p_{\text{corrected}} < 0.0001\)
    \item min\_prob: \(p_{\text{corrected}} = 0.0144\)
\end{itemize}

\subsubsection{Bootstrap Confidence Intervals}

\textbf{Purpose:} Bootstrap provides \textit{uncertainty quantification} for AUC estimates. Unlike parametric methods (e.g., DeLong's method), bootstrap makes minimal assumptions about the AUC's sampling distribution, making it robust when assumptions are violated or sample sizes are moderate.

\textbf{Value:} Confidence intervals allow us to distinguish between ``AUC significantly better than chance (0.5)'' versus ``AUC estimate with large uncertainty overlapping 0.5,'' providing evidence strength beyond point estimates alone.

\paragraph{Algorithm} (\(n_{\text{boot}} = 10{,}000\) resamples):
\begin{enumerate}
    \item For \(b = 1, \ldots, n_{\text{boot}}\):
    \begin{enumerate}
        \item Resample \(n\) edges with replacement: \((X^{(b)}, Y^{(b)})\)
        \item Compute \(\text{AUC}^{(b)}\) on bootstrap sample
    \end{enumerate}
    \item Compute percentile 95\% CI:
    \[
    \text{CI}_{95\%} = [\text{AUC}_{2.5\%}, \text{AUC}_{97.5\%}]
    \]
    where percentiles computed from empirical bootstrap distribution \(\{\text{AUC}^{(1)}, \ldots, \text{AUC}^{(n_{\text{boot}})}\}\)
\end{enumerate}

\textbf{Stratified resampling:} Maintains class proportions in each bootstrap sample to avoid degenerate cases with single class.

\textbf{Bootstrap standard error:}
\[
\text{SE}_{\text{boot}} = \sqrt{\frac{1}{n_{\text{boot}}-1}\sum_{b=1}^{n_{\text{boot}}}\left(\text{AUC}^{(b)} - \overline{\text{AUC}}\right)^2}
\]

\paragraph{Obtained Results.} Single-CLD analysis (Depressive Symptoms, n=182):
\begin{itemize}
    \item \textbf{max\_window\_entropy}: AUC = 0.762, 95\% CI [0.680, 0.836], \(\text{SE}_{\text{boot}} = 0.040\)
    \begin{itemize}
        \item CI excludes 0.5 \(\rightarrow\) significantly better than chance
        \item Narrow CI indicates stable, reliable estimate
    \end{itemize}
    \item \textbf{perplexity}: AUC = 0.747, 95\% CI [0.657, 0.829], \(\text{SE}_{\text{boot}} = 0.044\)
    \begin{itemize}
        \item CI excludes 0.5 \(\rightarrow\) significantly better than chance
    \end{itemize}
    \item \textbf{min\_prob}: AUC = 0.602, 95\% CI [0.505, 0.694], \(\text{SE}_{\text{boot}} = 0.048\)
    \begin{itemize}
        \item CI barely excludes 0.5, modest performance
    \end{itemize}
    \item \textbf{cosine\_similarity}: AUC = 0.224, 95\% CI [0.149, 0.305], \(\text{SE}_{\text{boot}} = 0.040\)
    \begin{itemize}
        \item CI entirely below 0.5 \(\rightarrow\) worse than chance (inverted predictor)
    \end{itemize}
\end{itemize}

Bootstrap distributions were approximately normal for all metrics (\texttt{rq2\_bootstrap\_ci.csv}, visualizations: \texttt{rq2\_bootstrap\_distributions.png}).

\subsubsection{Group Comparisons}

Independent samples t-tests compared CI metric distributions between hallucination and non-hallucination groups.

\paragraph{Test statistic:}
\[
t = \frac{\bar{x}_{\text{halluc}} - \bar{x}_{\text{non-halluc}}}{s_p \sqrt{\frac{1}{n_{\text{halluc}}} + \frac{1}{n_{\text{non-halluc}}}}}
\]

where \(s_p\) is pooled standard deviation:
\[
s_p = \sqrt{\frac{(n_{\text{halluc}}-1)s_{\text{halluc}}^2 + (n_{\text{non-halluc}}-1)s_{\text{non-halluc}}^2}{n_{\text{halluc}}+n_{\text{non-halluc}}-2}}
\]

\textbf{Degrees of freedom:} \(\nu = n_{\text{halluc}} + n_{\text{non-halluc}} - 2\)

\textbf{Hypothesis:} \(H_0: \mu_{\text{halluc}} = \mu_{\text{non-halluc}}\) vs. \(H_1: \mu_{\text{halluc}} \neq \mu_{\text{non-halluc}}\)

\paragraph{Effect Size.} Cohen's d quantifies magnitude of group differences standardized by pooled variance:
\[
d = \frac{\bar{x}_{\text{halluc}} - \bar{x}_{\text{non-halluc}}}{s_p}
\]

\textbf{Interpretation guidelines} \cite{cohen1988statistical}:
\begin{itemize}
    \item Small effect: \(|d| \approx 0.2\)
    \item Medium effect: \(|d| \approx 0.5\)
    \item Large effect: \(|d| \approx 0.8\)
\end{itemize}

\subsection{Meta-Analysis Framework}

\subsubsection{Cross-CLD Aggregation}

Meta-analytic statistics computed across \(n_{\text{CLD}} = 3\) independent datasets with 2 runs each.

\paragraph{Mean and Standard Deviation.} For statistic \(S\) (correlation \(r\) or AUC):
\[
\bar{S} = \frac{1}{n_{\text{analyses}}}\sum_{j=1}^{n_{\text{analyses}}} S_j, \quad s_S = \sqrt{\frac{1}{n_{\text{analyses}}-1}\sum_{j=1}^{n_{\text{analyses}}}(S_j - \bar{S})^2}
\]

\textbf{Standard deviation interpretation:} Quantifies cross-CLD stability. Low \(s_S\) (< 0.1) indicates consistent performance; high \(s_S\) (> 0.15) suggests domain-specific behavior.

\paragraph{Meta-Analytic Hypothesis Tests.}

\textbf{Correlation test:} \(H_0: \mu_r = 0\) (no relationship)
\[
t = \frac{\bar{r}}{s_r / \sqrt{n_{\text{analyses}}}}, \quad \nu = n_{\text{analyses}} - 1
\]

\textbf{AUC test:} \(H_0: \mu_{\text{AUC}} = 0.5\) (chance performance)
\[
t = \frac{\overline{\text{AUC}} - 0.5}{s_{\text{AUC}} / \sqrt{n_{\text{analyses}}}}, \quad \nu = n_{\text{analyses}} - 1
\]

\textbf{Implementation:} One-sample t-test comparing sample mean to null hypothesis value.

\paragraph{Confidence Intervals for Meta-Statistics.}
\[
\bar{S} \pm t_{\alpha/2, \nu} \cdot \frac{s_S}{\sqrt{n_{\text{analyses}}}}
\]

For 95\% CI with \(n_{\text{analyses}} = 32\), \(t_{0.025, 31} \approx 2.04\), but we use normal approximation \(z_{0.025} = 1.96\) for large samples:
\[
\text{CI}_{95\%} = \left[\bar{S} - 1.96\frac{s_S}{\sqrt{n_{\text{analyses}}}}, \bar{S} + 1.96\frac{s_S}{\sqrt{n_{\text{analyses}}}}\right]
\]

\paragraph{Obtained Meta-Analysis Results.} Aggregated across 32 analyses (3 CLDs \(\times\) 2 runs/CLD):

\textbf{Correlation meta-statistics} (testing \(H_0: \mu_r = 0\)):
\begin{itemize}
    \item \textbf{cosine\_similarity}: \(\bar{r} = +0.256\) [+0.232, +0.280], \(p < 0.001\)
    \item \textbf{max\_window\_entropy}: \(\bar{r} = +0.211\) [+0.175, +0.247], \(p < 0.001\)
    \item \textbf{perplexity}: \(\bar{r} = +0.200\) [+0.146, +0.254], \(p < 0.001\)
    \item \textbf{min\_prob}: \(\bar{r} = -0.162\) [-0.177, -0.147], \(p < 0.001\)
\end{itemize}

\textbf{AUC meta-statistics} (testing \(H_0: \mu_{\text{AUC}} = 0.5\)):
\begin{itemize}
    \item \textbf{perplexity}: \(\overline{\text{AUC}} = 0.679\) [0.650, 0.709], \(s = 0.076\), \(p < 0.001\)
    \begin{itemize}
        \item Consistently above chance, stable across domains
    \end{itemize}
    \item \textbf{max\_window\_entropy}: \(\overline{\text{AUC}} = 0.667\) [0.639, 0.695], \(s = 0.072\), \(p < 0.001\)
    \begin{itemize}
        \item Consistently above chance, stable across domains
    \end{itemize}
    \item \textbf{min\_prob}: \(\overline{\text{AUC}} = 0.641\) [0.617, 0.664], \(s = 0.060\), \(p < 0.001\)
    \begin{itemize}
        \item Low variance (\(s < 0.1\)) indicates high stability
    \end{itemize}
    \item \textbf{cosine\_similarity}: \(\overline{\text{AUC}} = 0.230\) [0.214, 0.245], \(s = 0.040\), \(p < 0.001\) vs.\ 0.5
    \begin{itemize}
        \item Significantly \textit{worse} than chance (inverted predictor)
    \end{itemize}
\end{itemize}

All top-performing metrics demonstrated statistically significant mean effects that were consistent across multiple independent CLDs (\texttt{rq2\_aggregate\_report\_*.md}, \texttt{aggregate\_meta\_statistics.xlsx}).

\subsection{Assumptions and Diagnostics}

\subsubsection{Parametric Test Assumptions}

\paragraph{Normality.} CI metrics exhibit right-skewed distributions (perplexity, max window entropy). However:
\begin{itemize}
    \item \textbf{Correlation tests:} Large sample sizes (\(n > 100\) per CLD) invoke Central Limit Theorem, relaxing normality requirements.
    \item \textbf{t-tests:} CLT justifies approximate normality of sampling distributions for group means.
    \item \textbf{Validation:} Spearman rank correlation provides non-parametric alternative; permutation tests require no distributional assumptions.
\end{itemize}

\paragraph{Homoscedasticity.} Equal variance assumption for t-tests checked via Levene's test. Violations addressed using Welch's t-test (unequal variances):
\[
t = \frac{\bar{x}_1 - \bar{x}_2}{\sqrt{s_1^2/n_1 + s_2^2/n_2}}, \quad \nu \approx \frac{(s_1^2/n_1 + s_2^2/n_2)^2}{(s_1^2/n_1)^2/(n_1-1) + (s_2^2/n_2)^2/(n_2-1)}
\]

\paragraph{Independence.} 
\begin{itemize}
    \item \textbf{Within-CLD:} Edges sharing variables may exhibit dependencies. Violates strict independence for correlation tests. Conservative interpretation: treat significance thresholds as approximate.
    \item \textbf{Between-CLD:} Independent datasets ensure valid meta-analytic aggregation.
\end{itemize}

\subsubsection{Multiple Testing Considerations}

Total hypothesis tests conducted:
\begin{itemize}
    \item 9 metrics \(\times\) 2 correlation types (Pearson, Spearman) = 18 correlation tests per CLD
    \item 9 metrics \(\times\) 1 AUC test per CLD
    \item 9 metrics \(\times\) 1 group comparison per CLD
\end{itemize}

\textbf{Correction strategy:}
\begin{itemize}
    \item \textbf{Primary analysis:} Holm-Bonferroni correction applied to 4 generator metrics (core hypotheses)
    \item \textbf{Exploratory analysis:} Judge metrics and aggregate score treated as exploratory; p-values reported without correction but interpreted cautiously
\end{itemize}

\textbf{Recommendation:} Results with \(p < 0.01\) robust to multiple testing; borderline results (\(0.01 < p < 0.05\)) require independent validation.

\subsection{Computational Implementation}

\subsubsection{Software Environment}

\begin{itemize}
    \item \textbf{Python:} 3.11.x
    \item \textbf{Core statistical libraries:} SciPy 1.11.x (statistical tests), scikit-learn 1.3.x (classification metrics), NumPy 1.25.x (numerical operations), pandas 2.0.x (data manipulation), statsmodels 0.14.x (multiple comparison correction)
\end{itemize}

\subsubsection{Analysis Pipeline Architecture}

Three-stage pipeline orchestrates statistical analyses:

\paragraph{Stage 1: Data Preparation.}
\begin{itemize}
    \item Loads experiment outputs and normalizes data format
    \item Defines binary hallucination labels based on ground truth or judge verdict
    \item Selects CI metric source (generator vs. judge metrics)
    \item Exports prepared data for downstream analysis
\end{itemize}

\paragraph{Stage 2: Single-CLD Analysis.}
\begin{itemize}
    \item \textbf{Correlation module:} Pearson and Spearman correlations with p-values
    \item \textbf{Classification module:} ROC curves, AUC computation, optimal threshold selection
    \item \textbf{Validation module:} Permutation tests (\(n=10{,}000\)), bootstrap CI (\(n=10{,}000\))
    \item \textbf{Comparison module:} Independent t-tests, Cohen's d effect sizes
    \item \textbf{Outputs:} CSV tables (correlations, AUC scores, thresholds, validation statistics)
\end{itemize}

\paragraph{Stage 3: Meta-Analysis Aggregation.}
\begin{itemize}
    \item Aggregates results across multiple CLDs
    \item Computes meta-statistics (mean, std, 95\% CI)
    \item One-sample t-tests: correlation \(\neq 0\), AUC > 0.5
    \item Identifies stable metrics (low cross-CLD variance)
    \item Generates aggregate visualizations (heatmaps, forest plots)
\end{itemize}

\subsubsection{Reproducibility Specifications}

\paragraph{Random seed control.} Stochastic procedures (permutation sampling, bootstrap resampling) employ explicit random seeds for reproducibility. Explicit seed values ensure exact replication across independent runs.

\paragraph{Computational cost.} Approximate runtime on standard hardware (Intel i7, 16GB RAM):
\begin{itemize}
    \item Single-CLD analysis (4 generator metrics, \(n_{\text{perm}} = n_{\text{boot}} = 10{,}000\)): 5--10 minutes
    \item Meta-analysis aggregation (3 CLDs, 32 analyses): 30--60 minutes total
    \item Memory footprint: < 2GB for largest CLD (\(N \approx 1100\) edges)
\end{itemize}

\paragraph{Parallelization.} Permutation and bootstrap loops executed serially. Potential speedup via \texttt{joblib.Parallel} or \texttt{multiprocessing.Pool} (not implemented in current version).

\subsection{Statistical Power and Sample Size}

\paragraph{Correlation detection power.} For two-tailed correlation test at \(\alpha = 0.05\), power to detect medium effect (\(r = 0.3\)) with 80\% power:
\[
n \approx 84 \text{ edges}
\]

All three CLDs exceed this minimum (Depressive symptoms: \(N=182\), Social norms: \(N=90\), Emergency visits: \(N=1119\)), providing adequate power for medium effects.

\paragraph{Meta-analysis power.} With 3 CLDs and multiple runs per CLD (\(n_{\text{analyses}} = 32\) for some metrics), meta-analytic tests achieve high power (> 90\%) for detecting consistent effects. However, small CLD count limits power for detecting heterogeneity (domain-specific effects).

\paragraph{Limitations.} Insufficient power for detecting small effects (\(r < 0.2\)) or weak classification performance (AUC < 0.6). Null findings should be interpreted as "insufficient evidence" rather than "evidence of absence."

\subsection{Summary}

This statistical framework combines classical parametric methods (correlation, t-tests) with robust non-parametric validation (permutation tests, bootstrap CI) and cross-dataset meta-analysis. Holm-Bonferroni correction controls family-wise error rate across primary hypotheses. Implementation in standard Python scientific stack ensures reproducibility and extensibility. Key methodological strengths include empirical null distributions (permutation testing), assumption-free uncertainty quantification (bootstrap), and hierarchical validation (single-CLD \(\rightarrow\) meta-analysis). Primary limitation: small CLD count (\(n=3\)) restricts power for heterogeneity analyses and domain stratification.
